\section{Core Idea}

Model the labyrinth as an unweighted graph on walkable cells.

We run two BFS passes:

\begin{itemize}
\item A \textit{multi-source BFS} from every monster cell computes
$monster\_dist[r][c]$, the earliest minute when any monster can reach cell
$(r,c)$.
\item A second BFS from \texttt{A} searches only cells that are still safe when
we arrive there.
\end{itemize}

The key invariant is the \textit{Safety Frontier Invariant}:

\textit{Every cell in the player's BFS queue is reached by a path whose arrival
time is strictly smaller than the earliest monster arrival time at every cell on
that path. In particular, for every queued cell $v$, we have
$player\_dist[v] < monster\_dist[v]$.}

That strict inequality is the whole problem. If we ever arrive at the same time
as the monster frontier, the cell is already unsafe.

So during the player's BFS, we only relax a neighbor $v$ when

$$
player\_dist[u] + 1 < monster\_dist[v].
$$

If we reach a boundary cell under this rule, then the reconstructed BFS path is
a valid escape. If no boundary cell is reachable, then no valid plan exists.

\section{Toy Example}

Use this tiny custom grid:

\begin{lstlisting}
..M..
.A...
.....
.....
\end{lstlisting}

Coordinates are zero-indexed. The player starts at $(1,1)$ and the monster
starts at $(0,2)$.

This example is intentionally small, but it exposes the decisive hard step:

\begin{itemize}
\item the top boundary cell $(0,1)$ looks attractive because it is only one move
away from \texttt{A}
\item however, the monster also reaches $(0,1)$ in one move, so that cell is
not safe
\item the left boundary cell $(1,0)$ is also one move away for the player, but
the monster needs three moves to reach it, so that exit is safe
\end{itemize}

The example is not about maze complexity. It is about why the condition must be
\textit{strictly earlier than the monster frontier}, not ``earlier or equal''.

\section{Walkthrough}

The decisive step is deciding which cells are allowed into the player's BFS
frontier. A full replay of both BFS passes would be long and low-value. The
hard part is the strict safety test.

In the toy example, the monster BFS gives these relevant arrival times:

\begin{itemize}
\item $monster\_dist(0,1) = 1$
\item $monster\_dist(1,2) = 1$
\item $monster\_dist(1,0) = 3$
\item $monster\_dist(2,1) = 3$
\end{itemize}

So from $(1,1)$, the top and right moves are rejected because the player would
arrive at exactly the same time as the monster frontier, while the left and
down moves stay safe.

\begin{animation}[id="monsters-safety-frontier", label="Monsters strict safety test"]
\shape{maze}{Grid}{rows=4, cols=5, data=[[".",".","M",".","."],[".","A",".",".","."],[".",".",".",".","."],[".",".",".",".","."]], label="toy grid"}
\shape{state}{VariableWatch}{names=["candidate","player_arrival","monster_arrival","safe"], label="Safety check"}

\step
\recolor{maze.cell[0][2]}{state=error}
\recolor{maze.cell[1][1]}{state=current}
\apply{state.var[candidate]}{value="start"}
\apply{state.var[player_arrival]}{value="0"}
\apply{state.var[monster_arrival]}{value="0 at M"}
\apply{state.var[safe]}{value="use monster BFS first"}
\narrate{First compute the earliest monster arrival time for every cell. Then run the player's BFS, but only through cells where the player arrives strictly earlier than the monsters.}

\step
\recolor{maze.cell[0][1]}{state=error}
\annotate{maze.cell[0][1]}{label="$1 \\not< 1$", arrow_from="maze.cell[1][1]", color=error}
\apply{state.var[candidate]}{value="(0,1)"}
\apply{state.var[player_arrival]}{value="1"}
\apply{state.var[monster_arrival]}{value="1"}
\apply{state.var[safe]}{value="no"}
\narrate{The top boundary cell looks attractive because it is one move away, but it fails the strict test. Arriving at the same time as the monster frontier is already unsafe.}

\step
\recolor{maze.cell[1][0]}{state=current}
\annotate{maze.cell[1][0]}{label="$1 < 3$", arrow_from="maze.cell[1][1]", color=good}
\apply{state.var[candidate]}{value="(1,0)"}
\apply{state.var[player_arrival]}{value="1"}
\apply{state.var[monster_arrival]}{value="3"}
\apply{state.var[safe]}{value="yes"}
\narrate{The left boundary cell survives. The player reaches it in one minute, while the earliest monster arrives only in three. This move stays inside the safe subgraph.}

\step
\recolor{maze.cell[1][1]}{state=path}
\recolor{maze.cell[1][0]}{state=path}
\apply{state.var[candidate]}{value="boundary"}
\apply{state.var[player_arrival]}{value="1"}
\apply{state.var[monster_arrival]}{value="3"}
\apply{state.var[safe]}{value="stop and reconstruct"}
\narrate{Once the player's BFS reaches a safe boundary cell, it can stop. Here the reconstructed escape path is just \texttt{L}. The second BFS is really BFS on the original grid after unsafe cells have been filtered out by the monster frontier.}
\end{animation}

\section{Proof Sketch}

There are two claims to prove.

\textbf{1. Monster BFS computes the correct danger time.}

Start the BFS queue with all monster cells at distance $0$. Because every move
costs one minute, standard multi-source BFS computes the shortest distance from
the set of monsters to every walkable cell. Therefore $monster\_dist[v]$ is the
earliest possible time when any monster can stand on $v$.

\textbf{2. Player BFS finds a valid escape if and only if one exists.}

We maintain the Safety Frontier Invariant:
every queued player cell $v$ satisfies
$player\_dist[v] < monster\_dist[v]$.

This is true for the start cell because \texttt{A} is not a monster cell.
Suppose it is true for the current cell $u$. We only push a neighbor $v$ when

$$
player\_dist[u] + 1 < monster\_dist[v].
$$

Then the path to $v$ is safe as well: every earlier cell on the path was already
safe by the invariant, and the new final cell $v$ is safe by the inequality
above. So the invariant is preserved.

This proves \textit{soundness}: every boundary cell reached by the player's BFS
gives a valid escape path.

For \textit{completeness}, consider any valid escape path. At every cell on that
path, the player must arrive strictly earlier than the earliest monster arrival;
otherwise a monster can be on that cell at the same time or sooner, which is
forbidden. So every valid escape path lies entirely inside the safe subgraph
defined by $player\_time < monster\_time$. The player's BFS explores exactly
that safe subgraph in increasing distance order. Therefore, if any valid escape
exists, the BFS reaches some boundary cell and reconstructs one.

\section{Complexity}

Both BFS passes touch each walkable cell at most once.

\begin{itemize}
\item Time complexity: $O(nm)$
\item Memory complexity: $O(nm)$
\end{itemize}

\section{Pitfalls}

\begin{itemize}
\item The safety test is strict. A move is allowed only when the player arrives
\textit{before} the monsters, so the condition is
\texttt{next\_dist < monster\_dist[nr][nc]}, not \texttt{<=}.
\item Monster distances must come from one multi-source BFS seeded with
\textit{all} monsters at once. Running separate BFS passes and taking minima is
slower and easier to get wrong.
\item Cells that monsters can never reach should keep distance \texttt{INF}. They
are still valid for the player if the player can reach them.
\item Path reconstruction needs the direction used to enter each cell, not just
the distance. Store a parent cell and a move letter when you push a neighbor.
\end{itemize}

\section{Reference Implementation}

\begin{lstlisting}[language=C++]
#include <algorithm>
#include <array>
#include <iostream>
#include <queue>
#include <string>
#include <utility>
#include <vector>

using namespace std;

int main() {
    int n, m;
    cin >> n >> m;

    vector<string> grid(n);
    for (int i = 0; i < n; ++i) {
        cin >> grid[i];
    }

    const int INF = 1e9;
    const array<int, 4> DR{1, -1, 0, 0};
    const array<int, 4> DC{0, 0, 1, -1};
    const array<char, 4> MOVE{'D', 'U', 'R', 'L'};

    auto inside = [&](int r, int c) {
        return 0 <= r && r < n && 0 <= c && c < m;
    };

    auto is_boundary = [&](int r, int c) {
        return r == 0 || r == n - 1 || c == 0 || c == m - 1;
    };

    vector<vector<int>> monster_dist(n, vector<int>(m, INF));
    vector<vector<int>> player_dist(n, vector<int>(m, -1));
    vector<vector<pair<int, int>>> parent(
        n, vector<pair<int, int>>(m, {-1, -1}));
    vector<vector<char>> parent_move(n, vector<char>(m, '?'));

    queue<pair<int, int>> q;
    pair<int, int> start{-1, -1};

    for (int r = 0; r < n; ++r) {
        for (int c = 0; c < m; ++c) {
            if (grid[r][c] == 'M') {
                monster_dist[r][c] = 0;
                q.push({r, c});
            } else if (grid[r][c] == 'A') {
                start = {r, c};
            }
        }
    }

    while (!q.empty()) {
        auto [r, c] = q.front();
        q.pop();

        for (int k = 0; k < 4; ++k) {
            int nr = r + DR[k];
            int nc = c + DC[k];
            if (!inside(nr, nc) || grid[nr][nc] == '#') {
                continue;
            }
            if (monster_dist[nr][nc] != INF) {
                continue;
            }
            monster_dist[nr][nc] = monster_dist[r][c] + 1;
            q.push({nr, nc});
        }
    }

    q.push(start);
    player_dist[start.first][start.second] = 0;

    pair<int, int> exit_cell{-1, -1};

    while (!q.empty()) {
        auto [r, c] = q.front();
        q.pop();

        if (is_boundary(r, c)) {
            exit_cell = {r, c};
            break;
        }

        for (int k = 0; k < 4; ++k) {
            int nr = r + DR[k];
            int nc = c + DC[k];
            if (!inside(nr, nc) || grid[nr][nc] == '#') {
                continue;
            }
            if (player_dist[nr][nc] != -1) {
                continue;
            }

            int next_dist = player_dist[r][c] + 1;
            if (next_dist >= monster_dist[nr][nc]) {
                continue;
            }

            player_dist[nr][nc] = next_dist;
            parent[nr][nc] = {r, c};
            parent_move[nr][nc] = MOVE[k];
            q.push({nr, nc});
        }
    }

    if (exit_cell.first == -1) {
        cout << "NO\n";
        return 0;
    }

    string path;
    for (pair<int, int> cur = exit_cell; cur != start; cur = parent[cur.first][cur.second]) {
        path += parent_move[cur.first][cur.second];
    }
    reverse(path.begin(), path.end());

    cout << "YES\n";
    cout << path.size() << '\n';
    cout << path << '\n';
    return 0;
}
\end{lstlisting}